\centerline{\bf \Large  РЕШЕНИЯ и КРИТЕРИИ}

\centerline{\bf \Large  7 КЛАСС}

\problem{\textbf{1}} Выпишем все возможные варианты трёх натуральных чисел, произведение которых равно 36.

$36\cdot1\cdot1=36$\\
$18\cdot2\cdot1=36$\\
$12\cdot3\cdot1=36$\\
$9\cdot4\cdot1=36$\\
$9\cdot2\cdot2=36$\\
$6\cdot6\cdot1=36$\\
$6\cdot3\cdot2=36$\\
$4\cdot3\cdot3=36$\\

Единственный вариант:

12, -3, -1.

Ответ: 12, -3, -1.
\par

\textbf{Критерии.} \\
\textit{7 баллов} $-$ Полностью верное решение\\
\textit{5 баллов} $-$ Правильный ответ без объяснений\\
\textit{0-3 балла} $-$  За все остальные попытки решить задачу\\

\problem{\textbf{2}} Пусть скорость Арслана $X$, скорость Батыра $Y$, а длина одного круга $S$.

Составим два уравнения:

$4X-4Y=9S$

$\dfrac{28}{3}X+\dfrac{28}{3}Y=24S$

Решив которое получим следующее равенство:

$X=15Y$

Ответ: в 15 раз быстрее.
\par
\textbf{Критерии.} \\
\textit{0-4 баллов} $-$ Решение пункта А\\
\textit{0-1 балла} $-$ Составление условия пункта Б \\
\textit{0-2 баллов} $-$ Решение пункта Б\\
\textit{0 баллов} $-$ Отсутствие решения пункта А (сразу НОЛЬ! баллов)\\
\textit{0 баллов} $-$ Любой ответ без объяснений\\


\problem{\textbf{3}} Оценка: $X$ не может быть меньше 31, иначе $X!$ не будет делиться на 31, так как 31 число простое.

Пример: $X!$ разделится на 18,19,20,...31 из определения факториала\par
$X!$ разделится на 32, так как $X!$ делится на 4 и на 8
$X!$ разделится на 33, так как $X!$ делится на 3 и на 11
$X!$ разделится на 34, так как $X!$ делится на 2 и на 17
$X!$ разделится на 35, так как $X!$ делится на 5 и на 7

\textbf{Критерии.} \\
\textit{7 баллов} $-$ Полностью верное решение\\
\textit{4 балла} $-$ Верный пример с объяснением делимости на все числа\\
\textit{3 балла} $-$ Верная оценка\\
\textit{0-1 балл} $-$ За все остальные попытки решить задачу\\
\textit{0 баллов} $-$ Любой ответ без объяснений\\


\problem{\textbf{4}} Найдём количество чисел, которые делятся на 5: 200/5=40

Количество чисел, которые делятся на 25: 200/25=8

Количество чисел, которые делятся на 5, но не делятся на 25: 40-8=32

Которые можем разбить на 16 пар и покрасить каждую в различные цвета.

Каждое число, которое делится на 25, покрасим также в различные цвета.

Максимум можем покрасить в 16+8=24 цвета все числа, чтобы условие задания выполнялось. Следовательно, Очир не может этого сделать.\par

\textbf{Критерии.} \\
\textit{7 баллов} $-$ Полностью верное решение\\
\textit{1 балл} $-$ Найдено количество чисел, которые делятся на 5\\
\textit{1 балл} $-$ Найдено количество чисел, которые делятся на 25\\
\textit{1 балл} $-$ Найдено количество чисел, которые делятся на 5, но не делятся на 25\\ 
\textit{1 балл} $-$ Разбиение 32-ух чисел на 16 пар и раскрашивание их в 16 различных цветов\\
\textit{1 балл} $-$ Раскрашивание чисел, которые делятся на 25 в 8 различных цветов\\
\textit{0-1 балл} $-$ За все остальные попытки решить задачу\\
\textit{0 баллов} $-$ Любой ответ без объяснений\\


\problem{\textbf{5}} Оценка: Отметим все числа с 1000000 до 10000000 с шагом в 1000. Таким образом, мы отметили все НЕУДЕкрасивые числа. Значит, УДЕкрасивых чисел может быть не большее 999, так как шаг в 1000.

Пример: Самое маленькое НЕУДЕкрасивое число - это $1000\cdot1000=1000000$, следующее же после него НЕУДЕкрасивое число - это $1000\cdot1001=1001000$. Между ними будут 999 УДЕкрасивых чисел.

1000001, 1000002, ... 1000999\par

\textbf{Критерии.} \\
\textit{7 баллов} $-$ Полностью верное решение\\
\textit{5 баллов} $-$ Верная оценка\\
\textit{2 балла} $-$ Верный пример с объяснением через минимальные НЕУДЕкрасивые числа\\
\textit{1 балл} $-$  Верный пример чисел с 1000001, 1000002, ... 1000999 без объяснения того, что между ними не может быть УДЕкрасивых\\ 
\textit{0-1 балл} $-$ За все остальные попытки решить задачу\\
